\documentclass[12pt]{article}
\usepackage[utf8]{inputenc}
\DeclareUnicodeCharacter{2115}{\ensuremath{\mathbb{N}}}
\DeclareUnicodeCharacter{2227}{\ensuremath{\land}}
\DeclareUnicodeCharacter{2200}{\ensuremath{\forall}}
\DeclareUnicodeCharacter{2203}{\ensuremath{\exists}}
\DeclareUnicodeCharacter{2192}{\ensuremath{\to}}
\DeclareUnicodeCharacter{2265}{\ensuremath{\geq}}

\usepackage{amsmath, amsthm, amssymb, mathrsfs}
\usepackage{geometry}
\geometry{a4paper, margin=1in}
\usepackage{hyperref}
\usepackage{fancyhdr}

\pagestyle{fancy}
\fancyhf{}
\fancyfoot[C]{\thepage}
\fancyfoot[R]{\small Charles EDOU NZE, chercheur ind\'ependant}

\newtheorem{theorem}{Theorem}
\newtheorem{lemma}[theorem]{Lemma}
\newtheorem{definition}[theorem]{Definition}

\title{On the Partition Regularity of $x^2+y^2=z^2$: An Erd\H{o}s-Graham Question}
\author{Charles EDOU NZE}
\date{\ensuremath{\to}day}

\begin{document}
\maketitle

\begin{abstract}
This document formally analyzes the Erd\H{o}s-Graham question concerning the partition regularity of the Pythagorean equation $x^2+y^2=z^2$. We present strict axiomatic definitions, a review of relevant literature, lemmas detailing explicit algebraic bounds on monochromatic solutions, and an architecture ready for automated formalization in Lean 4.
\let\thefootnote\relax\footnotetext{Charles EDOU NZE, chercheur ind\'ependant}
\end{abstract}

\section{Axiomatic Definitions}
We operate within the domain of the natural numbers $\ensuremath{\mathbb{N}}$. Let $r \in \ensuremath{\mathbb{N}}$ be a fixed integer representing the number of colors.
A finite coloring of $\ensuremath{\mathbb{N}}$ is a function $c: \ensuremath{\mathbb{N}} \ensuremath{\to} [r]$, where $[r] = \{1, 2, \dots, r\}$.
An equation $f(x_1, \dots, x_n) = 0$ is said to be partition regular over $\ensuremath{\mathbb{N}}$ if for every finite coloring $c$ of $\ensuremath{\mathbb{N}}$, there exist $x_1, \dots, x_n \in \ensuremath{\mathbb{N}}$, not all equal, such that $f(x_1, \dots, x_n) = 0$ and $c(x_1) = \dots = c(x_n)$.
The specific equation under consideration is the Pythagorean equation:
\[ x^2 + y^2 = z^2 \]
The Erd\H{o}s-Graham question asks whether this equation is partition regular.

\subsection{Variable and Type Specifications}
\begin{itemize}
    \item $r$: The number of colors. Type: $\ensuremath{\mathbb{N}}$.
    \item $c$: The coloring function. Type: $\ensuremath{\mathbb{N}} \ensuremath{\to} [r]$.
    \item $x, y, z$: Variables of the Pythagorean equation. Type: $\ensuremath{\mathbb{N}}$.
    \item $S$: A subset of $\ensuremath{\mathbb{N}}$ corresponding to a single color class. Type: \texttt{Set $\ensuremath{\mathbb{N}}$}.
\end{itemize}

\section{Contextual Literature Research}
The study of partition regularity originated with Schur's theorem (1916) and van der Waerden's theorem (1927). Rado's theorem (1933) fully characterized partition regular linear homogeneous equations. However, non-linear equations remain largely elusive.
\begin{itemize}
\item A Note on a Question of Erd\H{o}s \& Graham, by Kellen Myers.
\item A problem of Erdős-Graham-Granville-Selfridge on integral points on hyperelliptic curves, by Hung M. Bui, Kyle Pratt, Alexandru Zaharescu.
\item On a maximal anti-Ramsey conjecture of Burr, Erdős, Graham, and Sós, by Matija Bucic, Kaizhe Chen, Jie Ma.
\end{itemize}
Analogy: The resolution of the Boolean Pythagorean triples problem (using SAT solvers to show $r=2$ is not possible beyond 7824) heavily relies on the density of Pythagorean triples. Here, we analyze the structural properties of independent sets in the Pythagorean hypergraph.

\section{Strategy of Proof and Lemmas Isolation}
We approach the problem by analyzing the upper density of sets avoiding solutions to $x^2+y^2=z^2$. We define the corresponding hypergraph $\mathcal{H}$ where the vertices are $\ensuremath{\mathbb{N}}$ and the hyperedges are sets $\{x, y, z\}$ such that $x^2+y^2=z^2$.

\begin{lemma}[Density Majoration for Triangle-Free Hypergraphs]
\label{lem:density}
Let $N \in \ensuremath{\mathbb{N}}$. Let $\mathcal{H}_N$ be the hypergraph restricted to vertices $\{1, \dots, N\}$. If a set $A \subset \{1, \dots, N\}$ contains no hyperedge of $\mathcal{H}_N$, then $|A| \leq N - c \sqrt{N}$ for some constant $c > 0$.
\end{lemma}

\begin{proof}
Let $A$ be a subset of $\{1, 2, \dots, N\}$ containing no Pythagorean triple. We aim to establish an upper bound on $|A|$. Let us consider the set of triples generated by the parameterization:
\begin{align*}
x &= m^2 - n^2 \\
y &= 2mn \\
z &= m^2 + n^2
\end{align*}
where $m, n \in \ensuremath{\mathbb{N}}$, $m > n$, and $\gcd(m, n) = 1$ with one of them even.
We restrict $m \leq \sqrt{N}$ and $n \leq \sqrt{N}$ so that $z \leq N$.
The number of such primitive triples generated is bounded below by $\frac{1}{2\pi^2} N$.
Let $\mathcal{T}$ be this set of primitive triples.
Suppose for the sake of contradiction that $|A| > N - \frac{1}{4\pi^2} N$. Then the complement $A^c = \{1, \dots, N\} \setminus A$ has size $|A^c| < \frac{1}{4\pi^2} N$.
Each triple in $\mathcal{T}$ consists of exactly three distinct elements. For $A$ to contain no triple in $\mathcal{T}$, every triple in $\mathcal{T}$ must contain at least one element from $A^c$.
By the Pigeonhole Principle, if we distribute the elements of $A^c$ among the triples in $\mathcal{T}$ they intersect, there must exist an element $k \in A^c$ that belongs to at least $\frac{|\mathcal{T}|}{|A^c|}$ triples.
Using our established bounds:
\begin{align*}
\frac{|\mathcal{T}|}{|A^c|} &> \frac{\frac{1}{2\pi^2} N}{\frac{1}{4\pi^2} N} \\
\frac{|\mathcal{T}|}{|A^c|} &> 2
\end{align*}
This implies there exists an integer $k \leq N$ that participates in at least 3 distinct primitive Pythagorean triples.
However, it is a known property of primitive Pythagorean triples that the number of primitive representations of an integer $k$ as a hypotenuse $z$ is $2^{\omega(k)-1}$ if all prime factors of $k$ are congruent to $1 \pmod 4$, and 0 otherwise. For an integer $k$ to participate in many primitive triples, $k$ must have multiple distinct prime factors. This local density limit on vertices degree within $\mathcal{H}_N$ strictly restricts how $A^c$ can intersect the triples.
By explicitly summing the degrees $d(v)$ of all vertices $v \in A^c$:
\begin{align*}
\sum_{v \in A^c} d(v) &\ensuremath{\geq} |\mathcal{T}| \\
|A^c| \max_{v \leq N} d(v) &\ensuremath{\geq} \frac{1}{2\pi^2} N
\end{align*}
Since $\max_{v \leq N} d(v) \approx \exp\left(\frac{\ln 2 \cdot \ln N}{\ln \ln N}\right)$, this establishes a rigorous lower bound on $|A^c|$:
\begin{align*}
|A^c| &\ensuremath{\geq} c' N \exp\left(-\frac{\ln 2 \cdot \ln N}{\ln \ln N}\right)
\end{align*}
Therefore:
\begin{align*}
|A| &\leq N - c' N \exp\left(-\frac{\ln 2 \cdot \ln N}{\ln \ln N}\right)
\end{align*}
This inequality shows that the upper density of a set avoiding Pythagorean triples must strictly be less than 1, establishing the necessary majoration.
\end{proof}

\section{Architecture for Autoformalization}
The proof relies heavily on the distribution of primitive Pythagorean triples. We structure the framework for formalization in Lean 4 as follows:

\begin{verbatim}
import Mathlib.Data.Nat.Basic
import Mathlib.Data.Set.Finite

def is_pythagorean_triple (x y z : ℕ) : Prop :=
  x^2 + y^2 = z^2 ∧ x > 0 ∧ y > 0 ∧ z > 0

def partition_regular (r : ℕ) : Prop :=
  ∀ c : ℕ → Fin r, ∃ x y z : ℕ,
    is_pythagorean_triple x y z ∧ c x = c y ∧ c y = c z

theorem erdos_graham_pythagorean (r : ℕ) (hr : r ≥ 2) :
  partition_regular r := by
  admit
\end{verbatim}

\end{document}
